\section*{Capítulo \ref{ch:g11:polarity-and-cohesion} --- Electronegatividad, polaridad y fuerzas intermoleculares}

\begin{solution}{exo:g11:polarity-and-cohesion:1}
\ce{H-Cl}: el cloro (3.16 frente a 2.20). \ce{O-H}: el oxígeno (3.44).
\ce{C-O}: el oxígeno (3.44 frente a 2.55).
\end{solution}

\begin{solution}{exo:g11:polarity-and-cohesion:2}
\ce{H-H}: 0, no polar. \ce{C-H}: 0.35, se considera apolar. \ce{O-H}:
1.24, polar. \ce{N-H}: 0.84, polar. \ce{Cl-Cl}: 0, no polar. \ce{C-Cl}:
0.61, polar.
\end{solution}

\begin{solution}{exo:g11:polarity-and-cohesion:3}
El nitrógeno (3.04 > 2.19) y el oxígeno (3.44 > 2.58): la
\omterm{def:g11:polarity-and-cohesion:electronegativity}{electronegatividad} disminuye al bajar por un grupo. El nitrógeno
(3.04 > 2.55) y el magnesio (1.31 > 0.93): aumenta de izquierda a
derecha a lo largo de un periodo.
\end{solution}

\begin{solution}{exo:g11:polarity-and-cohesion:4}
El amoniaco y el agua: los dos tienen \omterm{def:g7:atoms-and-molecules:atom}{átomos} de hidrógeno unidos al
nitrógeno o al oxígeno, y \omterm{def:g10:lewis-and-shape:lone-pair}{pares no enlazantes} en ese \omterm{def:g7:atoms-and-molecules:atom}{átomo}. El metano no
tiene hidrógeno unido a N, O o F; el cloruro de hidrógeno tiene el
hidrógeno unido al cloro, que no es uno de los tres \omterm{def:g7:atoms-and-molecules:atom}{átomos} de la
definición.
\end{solution}

\begin{solution}{exo:g11:polarity-and-cohesion:5}
Sí: entre todas las \omterm{def:g7:atoms-and-molecules:molecule}{moléculas} hay \omterm{def:g11:polarity-and-cohesion:van-der-waals}{interacciones de van der Waals}. Pero
las \omterm{def:g7:atoms-and-molecules:molecule}{moléculas} de metano son pequeñas (10 \omterm{def:g9:inside-the-atom:nucleus}{electrones}), así que esas
interacciones son muy débiles y se rompen a muy baja temperatura: el
metano hierve hacia \qty{-162}{\celsius}.
\end{solution}

\begin{solution}{exo:g11:polarity-and-cohesion:6}
\ce{CH2Cl2}: polar (dos \omterm{def:g11:polarity-and-cohesion:polar-bond}{enlaces polares} \ce{C-Cl} y dos enlaces
\ce{C-H} casi apolares: ninguna simetría los anula). \ce{CCl4}: apolar
(cuatro enlaces idénticos en los vértices de un tetraedro). \ce{NH3}:
polar (pirámide). \ce{CO2}: apolar (lineal, enlaces opuestos).
\ce{BF3}: apolar (tres \omterm{def:g11:polarity-and-cohesion:polar-bond}{enlaces polares} idénticos a \ang{120} en un plano
se anulan).
\end{solution}

\begin{solution}{exo:g11:polarity-and-cohesion:7}
Se puede ceder el hidrógeno del grupo \ce{O-H} (no los del \ce{CH3},
unidos al carbono); el \omterm{def:g7:atoms-and-molecules:atom}{átomo} de oxígeno, con sus dos \omterm{def:g10:lewis-and-shape:lone-pair}{pares no
enlazantes}, puede recibir. Sí: el \ce{O-H} del metanol puede unirse al
oxígeno de una \omterm{def:g7:atoms-and-molecules:molecule}{molécula} de agua, y un \ce{O-H} del agua al oxígeno del
metanol, por ejemplo \ce{CH3-O-H}\,$\cdots$\,\ce{OH2}.
\end{solution}

\begin{solution}{exo:g11:polarity-and-cohesion:8}
Las tres \omterm{def:g7:atoms-and-molecules:molecule}{moléculas} son apolares y de la misma forma, pero cada vez más
grandes (10, 18 y 36 \omterm{def:g9:inside-the-atom:nucleus}{electrones}): sus \omterm{def:g11:polarity-and-cohesion:van-der-waals}{interacciones de van der Waals}
crecen, y con ellas sus temperaturas de ebullición.
\end{solution}

\begin{solution}{exo:g11:polarity-and-cohesion:9}
El \omterm{def:g9:periodic-table-first-look:period}{grupo} 14: el metano no tiene hidrógeno unido a N, O o F y no forma
\omterm{def:g11:polarity-and-cohesion:hydrogen-bond}{enlaces de hidrógeno}, así que sigue la tendencia de su grupo.
\end{solution}

\begin{solution}{exo:g11:polarity-and-cohesion:10}
Las \omterm{def:g11:polarity-and-cohesion:van-der-waals}{interacciones de van der Waals}: las \omterm{def:g7:atoms-and-molecules:molecule}{moléculas} crecen del \ce{HCl} al
\ce{HI}, y el aumento de esas interacciones con el tamaño pesa más que
la disminución de la atracción entre las \omterm{def:g11:polarity-and-cohesion:polar-bond}{cargas parciales}.
\end{solution}

\begin{solution}{exo:g11:polarity-and-cohesion:11}
Una \omterm{def:g7:atoms-and-molecules:molecule}{molécula} de yodo es mucho más grande que una de cloro (106
\omterm{def:g9:inside-the-atom:nucleus}{electrones} frente a 34): sus \omterm{def:g11:polarity-and-cohesion:van-der-waals}{interacciones de van der Waals} son mucho más
fuertes, lo bastante para mantener las \omterm{def:g7:atoms-and-molecules:molecule}{moléculas} en un sólido a
temperatura ambiente.
\end{solution}

\begin{solution}{exo:g11:polarity-and-cohesion:12}
El propano es apolar: solo \omterm{def:g11:polarity-and-cohesion:van-der-waals}{interacciones de van der Waals}. El
metoximetano es polar (dos \omterm{def:g11:polarity-and-cohesion:polar-bond}{enlaces polares} \ce{C-O} en un \ce{C-O-C}
angular), pero no tiene hidrógeno en su oxígeno: no hay \omterm{def:g11:polarity-and-cohesion:hydrogen-bond}{enlaces de
hidrógeno} entre sus \omterm{def:g7:atoms-and-molecules:molecule}{moléculas}. El etanol es polar y sus grupos \ce{O-H}
forman \omterm{def:g11:polarity-and-cohesion:hydrogen-bond}{enlaces de hidrógeno}. Cuanto más fuertes son las atracciones, más
alta es la temperatura de ebullición: propano < metoximetano < etanol.
\end{solution}

\begin{solution}{exo:g11:polarity-and-cohesion:13}
La red abierta del hielo ocupa más espacio que las mismas \omterm{def:g7:atoms-and-molecules:molecule}{moléculas} en
el líquido, así que el hielo es menos denso que el agua líquida y flota.
Un lago se hiela por arriba: la capa de hielo aísla el agua de debajo,
que sigue líquida, y los peces sobreviven.
\end{solution}

\begin{solution}{exo:g11:polarity-and-cohesion:14}
Metano sólido (\omterm{def:g11:polarity-and-cohesion:van-der-waals}{interacciones de van der Waals} entre \omterm{def:g7:atoms-and-molecules:molecule}{moléculas} pequeñas)
< hielo (\omterm{def:g11:polarity-and-cohesion:hydrogen-bond}{enlaces de hidrógeno}) < cloruro de potasio (atracción entre
\omterm{def:g9:ions:ion}{iones}).
\end{solution}

\begin{solution}{exo:g11:polarity-and-cohesion:15}
Una \omterm{def:g7:atoms-and-molecules:molecule}{molécula} de agua tiene dos \omterm{def:g7:atoms-and-molecules:atom}{átomos} de hidrógeno que ceder y dos \omterm{def:g10:lewis-and-shape:lone-pair}{pares
no enlazantes} con los que recibir: cada \omterm{def:g7:atoms-and-molecules:molecule}{molécula} puede participar en
cuatro \omterm{def:g11:polarity-and-cohesion:hydrogen-bond}{enlaces de hidrógeno}, dos cedidos y dos recibidos, lo que
construye una red completa. Una \omterm{def:g7:atoms-and-molecules:molecule}{molécula} de fluoruro de hidrógeno tiene
tres \omterm{def:g10:lewis-and-shape:lone-pair}{pares no enlazantes} pero un solo \omterm{def:g7:atoms-and-molecules:atom}{átomo} de hidrógeno: cede un solo
\omterm{def:g11:polarity-and-cohesion:hydrogen-bond}{enlace de hidrógeno}, así que en promedio solo se puede formar un enlace
por \omterm{def:g7:atoms-and-molecules:molecule}{molécula} (cadenas en zigzag). El agua tiene alrededor del doble de
\omterm{def:g11:polarity-and-cohesion:hydrogen-bond}{enlaces de hidrógeno} por \omterm{def:g7:atoms-and-molecules:molecule}{molécula}, y hierve más alto.
\end{solution}

\begin{solution}{pb:g11:polarity-and-cohesion:1}
\textbf{1.} \ce{H-O-H}, \ce{H-S-H}, \ce{H-Se-H}, cada \omterm{def:g7:atoms-and-molecules:atom}{átomo} central con
dos \omterm{def:g10:lewis-and-shape:lone-pair}{pares no enlazantes}: O, S y Se tienen seis \omterm{def:g10:electron-shells:valence}{electrones de valencia},
porque están en el mismo grupo.

\textbf{2.} $M(\ce{H2O}) = \qty{18.0}{g/mol}$, $M(\ce{H2S}) =
\qty{34.1}{g/mol}$, $M(\ce{H2Se}) = \qty{81.0}{g/mol}$.

\textbf{3.} $3.44 - 2.20 = 1.24$; $2.58 - 2.20 = 0.38$;
$2.55 - 2.20 = 0.35$. Solo el enlace \ce{O-H} es claramente polar.

\textbf{4.} Angular (cuatro grupos, dos de ellos \omterm{def:g10:lewis-and-shape:lone-pair}{pares no enlazantes}).
El agua es polar; el sulfuro de hidrógeno y el seleniuro de hidrógeno,
muy poco.

\textbf{5.} 10, 18 y 36 \omterm{def:g9:inside-the-atom:nucleus}{electrones}: el seleniuro de hidrógeno tiene las
\omterm{def:g11:polarity-and-cohesion:van-der-waals}{interacciones de van der Waals} más fuertes.

\textbf{6.} Solo el agua: sus \omterm{def:g7:atoms-and-molecules:atom}{átomos} de hidrógeno están unidos al
oxígeno, uno de los tres \omterm{def:g7:atoms-and-molecules:atom}{átomos} muy electronegativos; el azufre y el
selenio no son lo bastante electronegativos para dar a sus \omterm{def:g7:atoms-and-molecules:atom}{átomos} de
hidrógeno un $\delta^+$ marcado.

\textbf{7.} Cuatro: dos a través de sus propios \omterm{def:g7:atoms-and-molecules:atom}{átomos} de hidrógeno, y
dos recibidos por sus dos \omterm{def:g10:lewis-and-shape:lone-pair}{pares no enlazantes}.

\textbf{8.} \omterm{def:g11:polarity-and-cohesion:van-der-waals}{Interacciones de van der Waals} (incluida una débil
atracción entre las pequeñas \omterm{def:g11:polarity-and-cohesion:polar-bond}{cargas parciales}).

\textbf{9.} $212.9 - 273.15 = \qty{-60.3}{\celsius}$ y
\qty{-41.3}{\celsius}.

\textbf{10.} $-41.3 - (-60.3) = \qty{19.0}{\celsius}$.

\textbf{11.} La \omterm{def:g7:atoms-and-molecules:molecule}{molécula} es más grande, así que sus \omterm{def:g11:polarity-and-cohesion:van-der-waals}{interacciones de van
der Waals} son más fuertes.

\textbf{12.} $-60.3 - 19.0 = \qty{-79.3}{\celsius}$.

\textbf{13.} Subida: $-88.1 - (-112) = \qty{23.9}{\celsius}$; predicción
para el metano: $-112 - 23.9 = \qty{-136}{\celsius}$, frente a
\qty{-162}{\celsius} en la realidad: la extrapolación da unos
\qty{26}{\celsius} de más. No es exacta; la tendencia real no es una
recta.

\textbf{14.} Subida: $\qty{18.7}{\celsius}$; predicción para el fluoruro
de hidrógeno: $-85.1 - 18.7 = \qty{-103.8}{\celsius}$, frente a
\qty{19.6}{\celsius}: una diferencia de unos \qty{123}{\celsius}.

\textbf{15.} $100 - (-79.3) = \qty{179}{\celsius}$, unos
\qty{180}{\celsius}.

\textbf{16.} Subida: $\qty{25.3}{\celsius}$; predicción:
$-87.8 - 25.3 =
\qty{-113}{\celsius}$; diferencia:
$-33.4 - (-113) = \qty{80}{\celsius}$.

\textbf{17.} El agua (\qty{180}{\celsius}) > el fluoruro de hidrógeno
(\qty{123}{\celsius}) > el amoniaco (\qty{80}{\celsius}). El agua cede
dos \omterm{def:g7:atoms-and-molecules:atom}{átomos} de hidrógeno y recibe con dos \omterm{def:g10:lewis-and-shape:lone-pair}{pares no enlazantes}: cuatro
\omterm{def:g11:polarity-and-cohesion:hydrogen-bond}{enlaces de hidrógeno} por \omterm{def:g7:atoms-and-molecules:molecule}{molécula}, todos aprovechados. El fluoruro de
hidrógeno solo tiene un hidrógeno que ceder, y el amoniaco un solo \omterm{def:g10:lewis-and-shape:lone-pair}{par
no enlazante} con el que recibir: en promedio, menos enlaces por
\omterm{def:g7:atoms-and-molecules:molecule}{molécula}, y el nitrógeno, menos electronegativo, los forma más débiles.

\textbf{18.} Un gas: herviría hacia \qty{-80}{\celsius}. La Tierra no
tendría agua líquida, ni océanos, ni lluvia, ni nada de la vida que
depende de ellos.

\textbf{19.} Extrapola una recta trazada por solo dos puntos, y la prueba
con el \omterm{def:g9:periodic-table-first-look:period}{grupo} 14 muestra que una extrapolación así puede equivocarse en
unos \qty{25}{\celsius}.

\textbf{20.} Sin \omterm{def:g11:polarity-and-cohesion:hydrogen-bond}{enlaces de hidrógeno}, el agua herviría hacia
\qty{-80}{\celsius}, unos \qty{180}{\celsius} por debajo de su punto de
ebullición real.
\end{solution}
